\newpage
\section*{Введение}
\addcontentsline{toc}{section}{Введение}

В данной работе рассматриваются точки Лагранжа в ограниченной круговой задаче
трёх тел. Эта тема относится к небесной механике и имеет как теоретическое,
так и прикладное значение. С одной стороны, точки Лагранжа представляют собой
особые положения равновесия в системе двух массивных тел. С другой стороны,
они играют важную роль в задачах космонавтики, связанных с движением
космических аппаратов, размещением телескопов и исследованием устойчивых
областей в окрестности планет и спутников.

Обычная задача трёх тел является весьма сложной. В общем случае движение трёх
гравитирующих тел описывается системой нелинейных дифференциальных уравнений,
для которой не существует простого общего решения. Поэтому в практических
задачах часто рассматривается более простая модель - ограниченная круговая
задача трёх тел. В этой модели два тела большой массы движутся по круговым
орбитам вокруг общего центра масс, а третье тело имеет пренебрежимо малую
массу и не влияет на движение первых двух тел.

Наиболее интересным результатом этой модели является существование пяти
особых точек, в которых малое тело может находиться в состоянии относительного
равновесия. Эти точки называются точками Лагранжа или точками либрации.
Некоторые из них являются неустойчивыми, а некоторые при определённых условиях
обладают устойчивостью. Именно различие в поведении движения около таких точек
и представляет основной интерес в данной работе.

В данном проекте требуется не только разобрать теорию, связанную с точками
Лагранжа, но и запрограммировать движение малого тела в их окрестности.
Основная идея состоит в том, чтобы можно было задать начальную точку рядом с
любой из точек Лагранжа, численно проинтегрировать движение и получить
траекторию.

Особый интерес представляет сравнение устойчивых и неустойчивых точек. Если
задать начальные условия рядом с неустойчивой точкой Лагранжа, то траектория
должна со временем удаляться от неё. Если же начальные условия задать рядом с
устойчивой точкой, то движение должно оставаться в её окрестности достаточно
долго. Таким образом, вычислительная часть проекта должна наглядно подтвердить
теоретические выводы об устойчивости.

Следовательно, в работе ставятся две основные задачи. Первая задача -
рассмотреть математическую модель ограниченной круговой задачи трёх тел,
описать точки Лагранжа и объяснить, почему они бывают устойчивыми и
неустойчивыми. Вторая задача - реализовать вычислительный алгоритм, который
позволит находить точки Лагранжа, задавать начальные условия вблизи них и
строить траектории движения малого тела.

Итогом работы должна стать программа, позволяющая исследовать движение в
окрестности точек $L_1$-$L_5$, а также набор графиков, показывающих различие
между устойчивыми и неустойчивыми случаями. Тем самым проект объединяет
теоретическую часть небесной механики и её численную реализацию.
